% Geometric lemma for the exact signed roof-derivative continuation.
% To be inserted before the elementary pushforward lemma in native 336.
% The finite physical semialgebraic graph is inherited native 164
% (historical v74), not claimed as a newly discovered theorem.

\begin{lemma}[A common finite family of physical roof-side maps]
\label{lem:sbv-geometric-sides}
Fix the true collision count $m\ge1$. There is a finite integer $N_m$
and, for each original radius $R\in I$, a finite set
$\mathcal Z_{m,R}\subset[m\tau_-,m\tau_+]$, with
\[
                         \#\mathcal Z_{m,R}\le N_m,
\]
having the following property. In the signed integration-by-parts
formula for any original exact-record roof density in
Theorems~\ref{thm:rv-exact-bv} and \ref{thm:av-all-bv}, the regular
physical boundary terms can be written as a finite sum
\[
       (t_\gamma)_*(b_\gamma(r)\,dr).
\]
Here $r$ is a rational coordinate on a Gauss-direction arc,
$b_\gamma\in L^1$ is the signed or complex \emph{actual} one-copy
flux, and $t_\gamma$ is the original physical roof on that side.
Each $t_\gamma$ is $C^1$ and either constant, with value in
$\mathcal Z_{m,R}$, or has derivative nowhere zero on its open
parameter interval. The finitely many discarded parameter points
have zero flux. The positive-visited critical incoming contributions
are also supported on $\mathcal Z_{m,R}$.

The same set works simultaneously for every attainable exact label,
the marked phases, either $w=1$ or the original smooth guard
$w=W^\varepsilon$, and the bounded $C^1$ endpoint tests for which
Theorem~\ref{thm:wo-bv} proves finite variation. Neither a positive
incidence cutoff nor a replacement section is used. The set is
geometric: it is independent of $\varepsilon$, of the particular
smooth function $\Theta$, and of these test functions or phases.
No estimate uniform in $m$ is asserted for $N_m$.
\end{lemma}

\begin{proof}
We spell out the physical construction because local analyticity
alone would not give the asserted finite family uniformly in $R$.
The finite semialgebraic first-hit graph already proved in the
preserved native module 164 is the input here. We recall its
defining variables, rather than add a symbolic replacement model.

The uniform physical horizon gives a finite set of possible
next-center increments
\[
 \mathcal A=\{z\in\Lambda\setminus\{0\}:
                         |z|\le\tau_++2R_++1\}.
\]
For a word $C_0=0,\ldots,C_m$ with increments in $\mathcal A$,
use real contact points $Q_i=C_i+Rn_i$, unit normals $n_i$,
unit velocities $v_i$, and real lengths $\tau_i$.
The equations are
\[
 Q_{i+1}=Q_i+\tau_i v_i,\qquad
 v_i=v_{i-1}-2(v_{i-1}\cdot n_i)n_i,\qquad
 v_{i-1}\cdot n_i<0
\]
at every visited reflection, together with the original outgoing
sign at the initial collision and
$\tau_-\le\tau_i\le\tau_+$. To impose the first-hit rule, exclude
an interior intersection of the \emph{actual} flight with every
candidate disk preceding its designated endpoint. A disk that can
meet that segment lies within $\tau_++2R_+$ of its starting
center. Its squared distance to the segment is determined by the
two endpoint cases and by the interior orthogonal-foot case.
These are finite polynomial inequalities after the positive
projection denominator is retained. The current disk is treated
with its actual outgoing sign. No ray is restarted at a line foot.

The same finite enlargement gives every candidate for the original
least third-disk clearance.  A nearest-neighbor center other than
the terminal center has distance at most $1+R$ from the initial
contact, so its distance from the entire segment is at most $1+R$.
The least third-disk \emph{clearance} is therefore at most $1$.
A minimizing center is consequently within $\tau_++2R_++1$
of the starting center.
Square roots are encoded by their nonnegative root and squared
equation. Thus all physical incidences and original third-disk
clearances are semialgebraic on this graph.

Use finitely many rational charts for each normal circle, with
one-copy ownership and the finitely many seam sides retained until
cancellation. For a fixed chart rotation the charts have the form
\[
 n(x)=\operatorname{Rot}_{\alpha_*}
       \left(\frac{1-x^2}{1+x^2},\frac{2x}{1+x^2}\right),
 \qquad |x|<1.
\]
The moving center normal of each original rectangle is
semialgebraic in $R$, because its half-angle is the specified unique
root of the original polynomial in $q_R$. An angular side test is
$n\cdot n_*>\cos h$, with its actual fixed half-width $h$,
and its companion momentum test is linear. Fixed real numbers
such as $\cos h_0$ are allowed as coefficients.
Retain $\eta_0=\eta_m=1$,
$\sum_{i=0}^{m-1}\eta_i=n$, and $C_m=k$ exactly.
There are finitely many words and labels at this fixed $m$.
All these statements concern the deterministic physical graph and
the original roof $L=\sum_i\tau_i$.

On a regular endpoint sheet, write
$q=(\alpha_0,\alpha_m)$, $F=L$, and
$g=\nabla_qF=R(-p_0,p_m)$. The exact endpoint and optical
identities make $F$, $g$, and their derivatives semialgebraic in
rational normal charts: differentiation in $\alpha$ is
multiplication of the chart derivative by $(1+x^2)/2$.
Away from $g=0$, adjoin the positive variable $|g|$ and use
finitely many rational charts
\[
        e=\frac g{|g|}
           =\operatorname{Rot}_{\theta_*}
             \left(\frac{1-r^2}{1+r^2},
                         \frac{2r}{1+r^2}\right)
\]
for its true Gauss direction. The map $q\mapsto(F,e)$ has
the nonzero Jacobian $\chi$ of the retained optical coarea proof.
Consequently its regular fibers are discrete; since they are
semialgebraic, they are finite. The same statement holds for
the endpoint projection on a fixed regular word, by its nonzero
optical twist. Finite semialgebraic cell decomposition therefore
provides finitely many genuine inverse sheets, retaining their
physical multiplicities. It does not assign a nonphysical orbit
to a missing sheet.

Now take the complete one-sided graph of these sheets in variables
$(R,r,t)$, retaining the side identities before any observation
weight is inserted. Every closure and boundary in this construction
is taken \emph{fiberwise at the fixed radius}, before forming the
parameterized family. For example, membership of $(r,t)$ in the
closure of the $R$-fiber is expressed by
\[
 \forall a>0\ \exists(r',t')\quad
 (R,r',t')\text{ belongs to that sheet},\qquad
 (r-r')^2+(t-t')^2<a^2.
\]
Quantifier elimination makes this fiberwise closure semialgebraic.
The fiberwise relative boundary is obtained by the corresponding
Boolean operations. Thus limits from another radius never create
physical sides at an exceptional radius. At every fixed $R$ the
regular image is open in $(r,t)$, so its boundary has dimension at
most one. These boundaries include the original unvisited
constraints, caps used in the positive majorant, every actual
section side at every collision index, and their finite
intersections. They are semialgebraic, because closure, projection
and finite Boolean operations preserve semialgebraicity.
At positive visited incidences the critical set has
$g=0$ and positive endpoint Hessian. Its points are isolated on
each fixed-$R$ sheet. This is a zero-dimensional semialgebraic
fiber and hence finite. Over each such critical point the
Gauss-direction closure records its actual incoming sectors at
the single value $t=F(q)$. We keep those sectors, including
unvisited-tangency corners; they are not declared null because
the physical source point is null.

Apply $C^1$ cylindrical cell decomposition with $R$ ordered
before $r$. It is applied to the finite geometric side graphs,
their relative closures, and their projection-critical sets.
A side with nonconstant Gauss projection is then expressed on
finitely many cells as $t=t_\gamma(R,r)$.
For a fixed radius, further refine by the semialgebraic sets
where $\partial_rt_\gamma$ is zero, positive, or negative.
On an interval in the zero set the fundamental theorem of
calculus makes $t_\gamma$ constant. On each remaining interval
the derivative never vanishes. A side whose Gauss projection
is constant has zero $d\theta$-flux on its regular part and
does not create an omitted charge. Its intersections and the
finitely many point cells also have zero such flux.
The constant values and the critical roofs form
$\mathcal Z_{m,R}$.

This use of finite decomposition is precisely the finite
definable-family assertion in the cell and smooth-cell
decomposition theorems, for example
\cite[Theorems 2.9--2.10 and 6.6--6.7]{CosteOM}.
Their hypotheses apply to the semialgebraic graphs just
constructed. The finite cylindrical decomposition has finitely
many cells across the entire original radius interval, including
its exceptional parameter cells. Thus the number of its
fixed-radius interval pieces, and hence the number of their
constant values and critical roofs, has a finite common upper
bound $N_m$. These theorems provide no upper bound independent
of the formula size, which grows with $m$.

The inherited physical integration-by-parts formula on a roof
interval is the interior derivative plus its signed left and
right traces. For a regular side in the displayed chart,
$d\theta=2\,dr/(1+r^2)$ up to its chosen orientation.
Accordingly $b_\gamma(r)$ is the signed trace, including the
difference of the actual two incident record weights, times
this coordinate factor. Coincident genuine arcs have their
incident values combined before assigning the flux.
The physical trace and section-charge theorems give its $L^1$
bound; they do not require a bounded trace at each parameter
endpoint. Artificial angular charts, length cuts and
homogeneity cuts cancel between their two incident copies
before this representation. At a visited grazing limit the
optical trace is zero by Lemma~\ref{lem:rbv-endpoints},
including the simultaneous critical limit at its stated scope.
Thus it creates no additional side measure. The positive
visited critical fan is the separately retained incoming
measure described above.

Finally none of the smooth observation weights changes a
geometric sheet or introduces a jump side. The original
$W^\varepsilon$ is continuous on each physical sheet,
including where two branches of a least clearance exchange;
its interior derivative is the one already controlled in
the guarded BV proof. A bounded $C^1$ endpoint test similarly
adds only its proved interior derivative and multiplies
existing traces. Exact labels and discrete twists are constant
on the original section chambers. They therefore change
$b_\gamma$ and the critical sector coefficients but not the
finite side maps or their constant values. This proves the
claimed common set and its scope. An arbitrary discontinuous
Borel trajectory selector does not have this conclusion.
\end{proof}

% Bibliography item, to be added to the main bibliography if not already present:
% \bibitem{CosteOM} M. Coste, \emph{An Introduction to O-minimal Geometry},
% Istituti Editoriali e Poligrafici Internazionali, Pisa, 2000;
% author-hosted OMIN.pdf, Theorems 2.9--2.10 (p. 19) and 6.6--6.7 (p. 71).

